%%%%%%%%%%%%%%%%%%%%%%%%%%%%%%%%%%%%%%%%%%%%%%%%%%%%%%%%%%%%%%%%%%%%%%%%%%%%%%%
%  MA1521 Calculus for Computing -- Midterm Exercise Set
%  Typeset in the style of Thomas' Calculus (Early Transcendentals, 15/e),
%  laid out one-column in the manner of an MIT 18.02 problem set,
%  with working space after every problem.
%
%  Coverage: Thomas 11.1-11.5, 12.1-12.3, Chapter 13 (complete)
%            MIT 18.02 Lectures 1-13
%%%%%%%%%%%%%%%%%%%%%%%%%%%%%%%%%%%%%%%%%%%%%%%%%%%%%%%%%%%%%%%%%%%%%%%%%%%%%%%
\documentclass[11pt,oneside]{article}

\usepackage[T1]{fontenc}
\usepackage[utf8]{inputenc}
\usepackage{amsmath,amssymb,mathtools}
\usepackage[scale=0.92]{tgheros}
\usepackage{newtxtext}
\usepackage[varg]{newtxmath}
\usepackage[letterpaper,top=0.95in,bottom=0.95in,left=1.05in,right=1.05in,headsep=14pt,footskip=26pt]{geometry}
\usepackage{xcolor}
\usepackage{enumitem}
\usepackage{fancyhdr}
\usepackage{needspace}
\usepackage{lastpage}
\usepackage{graphicx}
\usepackage{multicol}
\usepackage{microtype}
\usepackage[hidelinks]{hyperref}

%------------------------------------------------------------------ palette --
\definecolor{ThomasBlue}{HTML}{00558C}   % Pearson/Thomas heading blue
\definecolor{ThomasDeep}{HTML}{003A62}
\definecolor{ThomasTint}{HTML}{E9F1F7}   % pale box fill
\definecolor{RuleGray}{HTML}{B9C4CC}
\definecolor{WorkGray}{HTML}{D9DEE3}
\definecolor{SoftGray}{HTML}{6A7378}

%------------------------------------------------------------------- vectors --
\renewcommand{\vec}[1]{\mathbf{#1}}
\newcommand{\bi}{\mathbf{i}}
\newcommand{\bj}{\mathbf{j}}
\newcommand{\bk}{\mathbf{k}}
\newcommand{\bA}{\mathbf{A}}
\newcommand{\bB}{\mathbf{B}}
\newcommand{\bC}{\mathbf{C}}
\newcommand{\bD}{\mathbf{D}}
\newcommand{\br}{\mathbf{r}}
\newcommand{\bv}{\mathbf{v}}
\newcommand{\bu}{\mathbf{u}}
\newcommand{\ba}{\mathbf{a}}
\newcommand{\bT}{\mathbf{T}}
\newcommand{\bn}{\mathbf{n}}
\newcommand{\grad}{\nabla}
\newcommand{\pd}[2]{\dfrac{\partial #1}{\partial #2}}
\newcommand{\pdi}[2]{\partial #1/\partial #2}
\newcommand{\dd}[2]{\dfrac{d#1}{d#2}}
\DeclareMathOperator{\proj}{proj}

%---------------------------------------------------------- running headers --
\pagestyle{fancy}
\fancyhf{}
\renewcommand{\headrulewidth}{0pt}
\renewcommand{\footrulewidth}{0pt}
\fancyhead[L]{\footnotesize\sffamily\color{SoftGray}MA1521 \textbar\ Midterm Exercise Set}
\fancyhead[R]{\footnotesize\sffamily\color{SoftGray}\leftmark}
\fancyfoot[C]{\footnotesize\sffamily\color{SoftGray}\thepage\ of \pageref{LastPage}}
\fancypagestyle{plain}{\fancyhf{}%
  \fancyfoot[C]{\footnotesize\sffamily\color{SoftGray}\thepage\ of \pageref{LastPage}}}

%--------------------------------------------------- Thomas-style structures --
% PART heading: heavy blue rule + "PART x" + title, in the manner of the
% "EXERCISES 13.7" banner that opens every Thomas exercise set.
\newcommand{\parthead}[4]{%
  \clearpage
  \markboth{#2}{#2}%
  \noindent{\color{ThomasBlue}\rule{\linewidth}{2.6pt}}\par\vspace{2pt}
  \noindent{\color{RuleGray}\rule{\linewidth}{0.5pt}}\par\vspace{7pt}
  \noindent{\sffamily\bfseries\fontsize{15}{17}\selectfont\color{ThomasDeep} PART #1}%
  \hfill{\sffamily\fontsize{9}{11}\selectfont\color{SoftGray} #4}\par\vspace{3pt}
  \noindent{\sffamily\bfseries\fontsize{21}{24}\selectfont\color{ThomasBlue} #2}\par\vspace{4pt}
  \noindent{\itshape\fontsize{10}{12}\selectfont\color{SoftGray} #3}\par\vspace{5pt}
  \noindent{\color{ThomasBlue}\rule{\linewidth}{1.1pt}}\par\vspace{12pt}
}

% Category sub-head, e.g. "Finding Local Extrema" in the textbook's exercise sets
\newcommand{\topichead}[1]{%
  \par\addvspace{15pt}\needspace{6\baselineskip}%
  \noindent{\sffamily\bfseries\fontsize{11.3}{13}\selectfont\color{ThomasBlue}#1}\par
  \vspace{1.5pt}\noindent{\color{RuleGray}\rule{\linewidth}{0.5pt}}\par\addvspace{7pt}
}

% Boxed DEFINITION / THEOREM panel, after the tinted panels in Thomas' text
\newcommand{\panel}[2]{%
  \par\addvspace{10pt}\needspace{5\baselineskip}%
  \noindent{\color{ThomasBlue}\rule{\linewidth}{1.1pt}}\par\vspace{-2pt}
  \noindent\colorbox{ThomasTint}{%
  \begin{minipage}{\dimexpr\linewidth-2\fboxsep\relax}
  \vspace{4pt}
  {\sffamily\bfseries\fontsize{9}{11}\selectfont\color{ThomasDeep}#1}\par\vspace{4pt}
  #2\par\vspace{4pt}
  \end{minipage}}\par\vspace{-1pt}
  \noindent{\color{RuleGray}\rule{\linewidth}{0.5pt}}\par\addvspace{11pt}
}

%------------------------------------------------------------------ problems --
\newcounter{prob}
\newcommand{\pts}[1]{{\sffamily\fontsize{8.5}{10}\selectfont\color{SoftGray}(#1)}}

% \begin{problem} ... \end{problem}
\newenvironment{problem}{%
  \par\addvspace{13pt}\needspace{5\baselineskip}%
  \refstepcounter{prob}%
  \noindent\hangindent=2.05em\hangafter=1
  {\bfseries\color{ThomasDeep}\theprob.}\hspace{0.55em}\ignorespaces
}{\par}

\newlist{parts}{enumerate}{3}
\setlist[parts,1]{label=\textbf{(\alph*)},ref=\alph*,leftmargin=2.15em,
                  itemsep=4pt,topsep=4pt,parsep=1pt,partopsep=0pt}
\setlist[parts,2]{label=(\roman*),leftmargin=2.0em,itemsep=3pt,topsep=3pt,parsep=0pt}

% Working space: a light ruled panel the student writes inside
\newlength{\wkwidth}
\newcommand{\work}[1]{%
  \par\addvspace{6pt}%
  \setlength{\wkwidth}{\dimexpr\linewidth-2\fboxsep-2\fboxrule\relax}%
  \noindent\fcolorbox{WorkGray}{white}{\rule{0pt}{#1}\rule{\wkwidth}{0pt}}%
  \par\addvspace{14pt}%
}

% Source attribution line
\newcommand{\src}[1]{\par\vspace{2pt}\hfill{\sffamily\fontsize{8}{9}\selectfont\color{SoftGray}[#1]}\par}

\setlength{\parindent}{0pt}
\setlength{\parskip}{0pt}
\linespread{1.04}

\begin{document}

%%%%%%%%%%%%%%%%%%%%%%%%%%%%%%%%%%%%%%%%%%%%%%%%%%%%%%%%%%%%%% TITLE BLOCK %%%
\thispagestyle{plain}
\markboth{Front Matter}{Front Matter}

\noindent{\color{ThomasBlue}\rule{\linewidth}{3.2pt}}\par\vspace{2.5pt}
\noindent{\color{RuleGray}\rule{\linewidth}{0.5pt}}\par\vspace{16pt}

\noindent
\begin{minipage}[t]{0.62\linewidth}
{\sffamily\bfseries\fontsize{12}{14}\selectfont\color{ThomasDeep}MA1521\ \ CALCULUS FOR COMPUTING}\par\vspace{2pt}
{\sffamily\fontsize{9.5}{11}\selectfont\color{SoftGray}National University of Singapore}
\end{minipage}%
\begin{minipage}[t]{0.38\linewidth}
\raggedleft
{\sffamily\fontsize{9.5}{12}\selectfont\color{SoftGray}
Exercise Set~M\ \ (Midterm Revision)\\
Name:\ \rule{3.2cm}{0.4pt}\\[2pt]
Date:\ \rule{3.2cm}{0.4pt}}
\end{minipage}

\vspace{20pt}
\noindent{\sffamily\bfseries\fontsize{28}{32}\selectfont\color{ThomasBlue}Vectors, Space Curves,\\[2pt] and Partial Derivatives}

\vspace{10pt}
\noindent{\fontsize{11.5}{14}\selectfont\itshape\color{ThomasDeep}
Thomas' Calculus: Early Transcendentals, 15/e \ \ \textbullet\ \ Sections 11.1--11.5,\ 12.1--12.3,\ and Chapter 13 complete}\par\vspace{3pt}
\noindent{\fontsize{10.5}{13}\selectfont\itshape\color{SoftGray}
MIT 18.02 Multivariable Calculus \ \ \textbullet\ \ Lectures 1--13}

\vspace{12pt}
\noindent{\color{ThomasBlue}\rule{\linewidth}{1.4pt}}
\vspace{16pt}

\panel{HOW TO USE THIS SET}{%
\fontsize{10}{12.5}\selectfont
The set is arranged in five parts. Parts~A and~B rebuild the geometric machinery
(\S11.1--11.5, \S12.1--12.3); Parts~C and~D are the heart of the midterm and cover
Chapter~13 in full, from level curves through Lagrange multipliers; Part~E is a
timed simulation assembled from all of the above.
\par\vspace{5pt}
Every problem is followed by a ruled panel sized to the work it should take. Do the
problem \emph{inside} the panel --- if your work spills out of the box, you are almost
certainly computing something the problem did not ask for. Answers to all problems
begin on the last pages; consult them only after a complete attempt.
\par\vspace{5pt}
Problems are drawn from the exercise sets of Thomas' Calculus 15/e, from the MIT 18.02
problem sets used in Lecture Group~1, and from the tutorial questions of Lecture Group~2.
The source of each part is recorded in its banner.
}

\vspace{4pt}

%%%%%%%%%%%%%%%%%%%%%%%%%%%%%%%%%%%%%%%%%%%%%%%%%%%%%%%%%%%%%%%%%%%%%%%%%%%%%%
\topichead{Distribution of Problems}

\noindent
{\fontsize{10}{13}\selectfont
\begin{tabular}{@{}p{0.06\linewidth}p{0.44\linewidth}p{0.26\linewidth}p{0.14\linewidth}@{}}
\textbf{Part} & \textbf{Topic} & \textbf{Thomas sections} & \textbf{Problems}\\[3pt]
\hline\\[-7pt]
A & Vectors and the Geometry of Space & 11.1--11.5 & 1--9\\
B & Vector-Valued Functions and Motion & 12.1--12.3 & 10--15\\
C & Partial Derivatives, Chain Rule, Gradient & 13.1--13.6 & 16--38\\
D & Extreme Values and Lagrange Multipliers & 13.7--13.9 & 39--52\\
E & Midterm Simulation (90 minutes) & 11--13 & 53--58\\[3pt]
\hline
\end{tabular}}

\vspace{4pt}

\topichead{Results You Are Expected to Know Cold}

\panel{THEOREM 1 --- Angle Between Two Vectors}{%
The angle $\theta$ between two nonzero vectors $\bu=\langle u_1,u_2,u_3\rangle$ and
$\bv=\langle v_1,v_2,v_3\rangle$ is
\[
\theta=\cos^{-1}\!\left(\frac{\bu\cdot\bv}{|\bu|\,|\bv|}\right),
\qquad
\bu\cdot\bv=u_1v_1+u_2v_2+u_3v_3 .
\]
\vspace{-9pt}
\[
\proj_{\bv}\bu=\left(\frac{\bu\cdot\bv}{|\bv|^{2}}\right)\bv,
\qquad
|\bu\times\bv|=|\bu|\,|\bv|\sin\theta,
\qquad
|\,(\bu\times\bv)\cdot\mathbf{w}\,|=\bigl|\det[\bu\ \bv\ \mathbf{w}]\bigr| .
\]}

\panel{DEFINITION --- Arc Length Along a Space Curve}{%
For a smooth curve $\br(t)=f(t)\bi+g(t)\bj+h(t)\bk$, $a\le t\le b$,
\[
L=\int_a^b|\bv(t)|\,dt=\int_a^b\sqrt{\left(\dd{x}{t}\right)^{2}+\left(\dd{y}{t}\right)^{2}+\left(\dd{z}{t}\right)^{2}}\,dt,
\qquad
\bT=\frac{\bv}{|\bv|}.
\]}

\panel{THEOREM --- The Chain Rule, the Gradient, and the Second Derivative Test}{%
If $w=f(x,y,z)$ is differentiable and $x,y,z$ are differentiable functions of $t$, then
$\dd{w}{t}=f_x\dd{x}{t}+f_y\dd{y}{t}+f_z\dd{z}{t}$, and
$df=f_x\,dx+f_y\,dy+f_z\,dz$.
\par\vspace{4pt}
The directional derivative of $f$ at $P_0$ in the direction of a \emph{unit} vector $\bu$ is
$D_{\bu}f=\grad f\big|_{P_0}\!\cdot\bu$; it is largest ($=|\grad f|$) in the direction of
$\grad f$ and zero along the level curve through $P_0$.
\par\vspace{4pt}
At a critical point $(a,b)$ of $f$, set $D=f_{xx}f_{yy}-f_{xy}^{\,2}$ evaluated at $(a,b)$. Then
$f$ has a local maximum if $D>0$ and $f_{xx}<0$; a local minimum if $D>0$ and $f_{xx}>0$;
a saddle point if $D<0$; and the test is inconclusive if $D=0$.
\par\vspace{4pt}
\emph{Lagrange multipliers.} To find the extreme values of $f$ subject to $g=0$,
solve $\grad f=\lambda\grad g$ together with $g=0$, then compare the values of $f$
at the solutions.}


%%%%%%%%%%%%%%%%%%%%%%%%%%%%%%%%%%%%%%%%%%%%%%%%%%%%%%%%%%%%%%%%%%% PART A %%%
\parthead{A}{Vectors and the Geometry of Space}%
{Thomas 11.1--11.5 \ \textbullet\ \ MIT 18.02 Lectures 1--5}%
{9 problems \ \textbar\ \ 25 points}

\topichead{Vectors, Lengths, and Directions}

\begin{problem}
\textbf{Unit vectors and displacement.}
\begin{parts}
\item For what value or values of $c$ will $\tfrac{1}{2}\bi-c\bj+c\bk$ be a unit vector?
\item A vector $\bA$ has magnitude $6$ and direction $(\bi+2\bj-2\bk)/3$. If its tail is at
      $(-2,0,1)$, where is its head?
\end{parts}
\src{MIT 18.02, Problem Set 1}
\end{problem}
\work{3.6cm}

\begin{problem}
Let $P$ and $Q$ be two points in space, and let $X$ be the midpoint of the line segment $PQ$.
Let $O$ be an arbitrary fixed point.
\begin{parts}
\item Show that, as vectors, $\overrightarrow{OX}=\tfrac{1}{2}\bigl(\overrightarrow{OP}+\overrightarrow{OQ}\bigr)$.
\item With the notation of part (a), assume instead that $X$ divides the segment $PQ$ in the
      ratio $r:s$, where $r+s=1$. Derive an expression for $\overrightarrow{OX}$ in terms of
      $\overrightarrow{OP}$ and $\overrightarrow{OQ}$.
\end{parts}
\src{MIT 18.02, Problem Set 1}
\end{problem}
\work{4.8cm}

\topichead{The Dot Product, Projections, and Angles}

\begin{problem}
\textbf{Angles between vectors.}
\begin{parts}
\item Find the angle between $\bi-\bk$ and $4\bi+4\bj-2\bk$, and the angle between
      $\bi+\bj+2\bk$ and $2\bi-\bj+\bk$.
\item Tell for what values of $c$ the vectors $c\bi+2\bj-\bk$ and $\bi-\bj+2\bk$ will
      (i) be orthogonal; (ii) form an acute angle.
\end{parts}
\src{MIT 18.02, Problem Set 1}
\end{problem}
\work{4.2cm}

\begin{problem}
Consider the vectors $\bA=6\bj+6\bk$ and $\bB=2\bi-\bj-\bk$. Write $\bA=\bC+\bD$ as a sum of
two vectors such that $\bC$ is parallel to $\bB$ and $\bD$ is perpendicular to $\bB$.
\par\vspace{3pt}
{\itshape\small Suggestion. First write $\bC=\alpha\bB$, so that $\bA=\alpha\bB+\bD$ for some
real number $\alpha$. Determine $\alpha$ using the dot product, then proceed to get
$\bC$ and $\bD$.}
\src{MA1521 L2, Q04.4}
\end{problem}
\work{4.4cm}

\topichead{Determinants and the Cross Product}

\begin{problem}
\textbf{Determinants and the cross product.}
\begin{parts}
\item Calculate
\[
\begin{vmatrix} -1 & 0 & 4\\ 1 & 2 & 2\\ 3 & -2 & -1\end{vmatrix}
\]
using the Laplace expansion by the cofactors of (i) the first row; (ii) the first column.
\item Find $\bA\times\bB$ if
      (i) $\bA=\bi-2\bj+\bk$, $\bB=2\bi-\bj-\bk$;\quad
      (ii) $\bA=2\bi-3\bk$, $\bB=\bi+\bj-\bk$.
\end{parts}
\src{MIT 18.02, Problem Set 1}
\end{problem}
\work{5.1cm}

\begin{problem}
Consider the points $A=(5,2,0)$, $B=(2,6,1)$, and $C=(2,4,7)$ in space.
\begin{parts}
\item Find the area of the parallelogram with two adjacent sides given by the line segments
      $AB$ and $AC$. Hence find the area of triangle $ABC$.
\item Find the volume of the parallelepiped spanned by $\bi+3\bk$, $2\bi+\bj-2\bk$,
      and $5\bi+4\bk$.
\item Suppose $A$, $B$, $C$, $D$ are four points in space with
      $\det\bigl(\overrightarrow{AB},\overrightarrow{AC},\overrightarrow{AD}\bigr)=0$.
      Must the four points be coplanar? Give a reason.
\end{parts}
\src{MA1521 L2, Q04.8--Q04.10}
\end{problem}
\work{5.4cm}

\topichead{Lines and Planes in Space}

\begin{problem}
Find equations for the following planes.
\begin{parts}
\item Through $(2,0,-1)$ and perpendicular to $\bi+2\bj-2\bk$.
\item Through the origin, $(1,1,0)$, and $(2,-1,3)$.
\item Through the points on the $x$-, $y$-, and $z$-axes where $x=a$, $y=b$, $z=c$
      respectively (all nonzero).
\item Through $(1,0,1)$ and $(0,1,1)$ and parallel to $\bi-\bj+2\bk$.
\end{parts}
\src{MIT 18.02, Problem Set 1}
\end{problem}
\work{6.0cm}

\begin{problem}
Find, in parametric form, equations for
\begin{parts}
\item the line through $(1,0,-1)$ and parallel to $2\bi-\bj+3\bk$;
\item the line through $(2,-1,-1)$ and perpendicular to the plane $x-y+2z=3$;
\item all lines passing through $(1,1,1)$ and lying in the plane $x+2y-z=2$;
\item the line through $P(4,2,1)$ and perpendicular to both of the lines
      \[
      \ell_1:\ x=3t+2,\quad y=-4t-1,\quad z=4t-9,
      \qquad
      \ell_2:\ x=3t-4,\quad y=-t+6,\quad z=5t .
      \]
\end{parts}
\src{MIT 18.02, Problem Set 1; MA1521 L2, Q05.2}
\end{problem}
\work{6.3cm}

\begin{problem}
\textbf{Lines meeting planes.}
\begin{parts}
\item Where does the line through $(0,1,2)$ and $(2,0,3)$ intersect the plane $x+4y+z=4$?
\item The line passing through $(1,1,-1)$ and perpendicular to the plane $x+2y-z=3$
      intersects the plane $2x-y+z=1$ at what point?
\item Find an equation of the plane that contains the point $(-1,2,1)$ and the line of
      intersection of the planes $x+y-z=2$ and $2x-y+3z=1$.
\end{parts}
\src{MIT 18.02, Problem Set 1; MA1521 L2, Q05.5}
\end{problem}
\work{6.5cm}


%%%%%%%%%%%%%%%%%%%%%%%%%%%%%%%%%%%%%%%%%%%%%%%%%%%%%%%%%%%%%%%%%%% PART B %%%
\parthead{B}{Vector-Valued Functions and Motion in Space}%
{Thomas 12.1--12.3 \ \textbullet\ \ MIT 18.02 Lectures 5--6}%
{6 problems \ \textbar\ \ 20 points}

\topichead{Curves in Space and Their Tangents}

\begin{problem}
For each of the following vector functions of time, calculate the velocity, the speed, the
unit tangent vector (in the direction of the velocity), and the acceleration.
\begin{parts}
\item $\br=e^{t}\bi+e^{-t}\bj$
\item $\br=t^{2}\bi+t^{3}\bj$
\item $\br=(1-2t^{2})\bi+t^{2}\bj+(-2+2t^{2})\bk$
\end{parts}
\src{MIT 18.02, Problem Set 2}
\end{problem}
\work{5.6cm}

\begin{problem}
Consider the vector function
$\bv(t)=(3t^{2}+2t)\bi+\pi\sin(\pi t)\bj+e^{t}\bk$.
\begin{parts}
\item Evaluate $\bv'(t)$, and find $\bv'(2)$.
\item Evaluate the indefinite integral $\displaystyle\int\bv(t)\,dt$.
\item Evaluate the definite integral $\displaystyle\int_0^1\bv(t)\,dt$.
\end{parts}
\src{MA1521 L2, Q05.10}
\end{problem}
\work{5.1cm}

\begin{problem}
Consider the curve in $\mathbb{R}^{3}$ given by
\[
\br(t)=(\ln t)\,\bi+\frac{t-1}{t+3}\,\bj+(t\ln t)\,\bk .
\]
Find parametric equations of the tangent line to the curve at the point where $t=1$.
\src{MA1521 L2, Q05.11}
\end{problem}
\work{4.6cm}

\topichead{Integrals of Vector Functions; Motion}

\begin{problem}
A particle moves in $\mathbb{R}^{3}$ so that its velocity at time $t$ is
\[
\bv(t)=(6t^{2}+4t)\,\bi+(8t^{3}+2t)\,\bj+\frac{1}{1+t^{2}}\,\bk .
\]
If the position of the particle at time $t=0$ is $\bi+2\bj+\pi\bk$, find the position of the
particle at time $t=1$.
\src{MA1521 L2, Q05.12}
\end{problem}
\work{4.6cm}

\begin{problem}
Two particles move in $\mathbb{R}^{3}$ so that their positions at time $t$ are
\[
\br_1(t)=t\,\bi+t^{2}\,\bj+t^{3}\,\bk ,
\qquad
\br_2(t)=(1+2t)\,\bi+(1+6t)\,\bj+(1+14t)\,\bk .
\]
\begin{parts}
\item Do the particles collide? Justify your answer.
\item Do their paths intersect? Justify your answer.
\end{parts}
\src{MA1521 L2, Q05.13}
\end{problem}
\work{5.6cm}

\topichead{Arc Length in Space}

\begin{problem}
\textbf{Lengths of curves in space.}
\begin{parts}
\item Find the unit tangent vector of
      $\br(t)=(2\cos t)\bi+(2\sin t)\bj+\sqrt{5}\,t\,\bk$, and the length of the portion
      $0\le t\le\pi$.
\item Find the arc length of the curve
      \[
      x=2t-2\sin t,\qquad y=2-\sqrt{2}\cos t,\qquad z=\sqrt{2}\cos t,
      \qquad 0\le t\le 4\pi .
      \]
\item Find the length of the curve $\br(t)=2t\,\bi+2t\,\bj+(1-t^{2})\,\bk$ from
      $(0,0,1)$ to $(2,2,0)$.
\end{parts}
\src{Thomas 12.3, Exercises 1 and 15; MA1521 L2, Q05.14}
\end{problem}
\work{7.1cm}


%%%%%%%%%%%%%%%%%%%%%%%%%%%%%%%%%%%%%%%%%%%%%%%%%%%%%%%%%%%%%%%%%%% PART C %%%
\parthead{C}{Partial Derivatives, the Chain Rule, and the Gradient}%
{Thomas 13.1--13.6 \ \textbullet\ \ MIT 18.02 Lectures 8--12}%
{23 problems \ \textbar\ \ 60 points}

\topichead{Functions of Several Variables; Limits and Continuity}

\begin{problem}
\textbf{Values, domains, and level sets.}
\begin{parts}
\item Let $\displaystyle f(x,y,z)=\frac{x-y}{y^{2}+z^{2}}$. Evaluate $f(3,-1,2)$ and
      $f\!\left(0,-\tfrac{1}{3},0\right)$.
\item Find and sketch the natural domain of
      $\displaystyle f(x,y)=\frac{1}{\ln\left(9-x^{2}-y^{2}\right)}$.
\item Describe geometrically the level sets of
      $f(x,y,z)=\sqrt{16-x^{2}-y^{2}-z^{2}}$.
\end{parts}
\src{MA1521 L2, Q06.1, Q06.2, Q06.4}
\end{problem}
\work{5.4cm}

\begin{problem}
\textbf{Level curves and graphs.}
\begin{parts}
\item Sketch five level curves for each of the following functions, and sketch the portion
      of the graph lying in the first octant:
      \[
      \text{(i) } 1-x-y,\qquad
      \text{(ii) } \sqrt{x^{2}+y^{2}},\qquad
      \text{(iii) } x^{2}+y^{2},\qquad
      \text{(iv) } 1-x^{2}-y^{2}.
      \]
\item Sketch the level curves $f(x,y)=k$ for $k=-2,0,1,2$, where
      $\displaystyle f(x,y)=\frac{x+y}{x-y}$.
\end{parts}
\src{MIT 18.02, Problem Set 3; MA1521 L2, Q06.3}
\end{problem}
\work{7.3cm}

\begin{problem}
Evaluate each limit, or show that it does not exist.
\begin{parts}
\item $\displaystyle\lim_{(x,y)\to(0,\pi/4)}\sec x\,\tan(x+y)$
\item $\displaystyle\lim_{(x,y)\to(2,0)}\frac{\sqrt{2x-y}-2}{2x-y-4}$
      \quad{\itshape\small(The identity $a^{2}-b^{2}=(a+b)(a-b)$ is useful.)}
\item $\displaystyle\lim_{(x,y)\to(0,0)}\frac{y}{x^{2}}$
\item Let $f(x,y)=\sqrt{x^{2}+y^{2}}$ for $(x,y)\neq(0,0)$ and $f(0,0)=1$.
      Find $\lim_{(x,y)\to(0,0)}f(x,y)$, and decide whether $f$ is continuous at $(0,0)$.
\end{parts}
\src{MA1521 L2, Q06.5, Q06.6, Q07.1}
\end{problem}
\work{6.6cm}

\topichead{Partial Derivatives}

\begin{problem}
Calculate the first partial derivatives of each of the following functions.
\begin{parts}
\item $4x^{3}-3x^{2}y^{2}+2x$
\item $\sin(3x+2y)$
\item $y\ln(4x+3y)$; find also the values of the first partials at $(2,1)$.
\item $x^{2}z-2yz^{3}$
\item $y^{2}z+x^{3}yz^{3}$; find also the values of the first partials at $(1,2,4)$.
\end{parts}
\src{MIT 18.02, Problem Set 3; MA1521 L2, Q07.2}
\end{problem}
\work{6.1cm}

\begin{problem}
\textbf{Mixed second partials.}
\begin{parts}
\item Verify that $f_{xy}=f_{yx}$ for
      (i) $x^{m}y^{n}$ ($m,n$ positive integers);\quad
      (ii) $\dfrac{x}{1+y}$;\quad
      (iii) $\cos(x^{2}+y)$.
\item Using $f_{xy}=f_{yx}$, tell for what value of the constant $a$ there exists a function
      $f(x,y)$ with $f_x=axy+3y^{2}$ and $f_y=x^{2}+6xy$; then find such a function by
      inspection.
\item Suppose $f(x,y)$ satisfies $f_x=3x^{2}+Axy$ and $f_y=3x^{2}+6y^{2}$, where $A$ is a
      constant. Using Clairaut's Theorem, determine $A$, and then find such a function
      $f(x,y)$. {\itshape\small(Such a function is not unique.)}
\end{parts}
\src{MIT 18.02, Problem Set 3; MA1521 L2, Q07.6}
\end{problem}
\work{6.8cm}

\begin{problem}
\textbf{Laplace's equation and implicit partials.}
\begin{parts}
\item Show that each of the following functions $w=f(x,y)$ satisfies the two-dimensional
      Laplace equation $w_{xx}+w_{yy}=0$:
      (i) $w=e^{ax}\sin ay$ ($a$ constant);\quad (ii) $w=\ln\!\left(x^{2}+y^{2}\right)$.
\item Using implicit differentiation, find $\pdi{z}{x}$ and $\pdi{z}{y}$ if
      \[
      x^{3}+2y^{3}+6xz^{3}-2=59z .
      \]
      {\itshape\small Here $z$ is regarded as a function of the independent variables $x$ and $y$,
      defined implicitly by the equation.}
\end{parts}
\src{MIT 18.02, Problem Set 3; MA1521 L2, Q07.3}
\end{problem}
\work{6.5cm}

\topichead{Tangent Planes and Normal Lines}

\begin{problem}
Give an equation of the tangent plane to each of the following surfaces at the point
indicated.
\begin{parts}
\item $z=xy^{2}$ at $(1,1,1)$
\item $w=y^{2}/x$ at $(1,2,4)$
\item the graph of $f(x,y)=x^{2}y+2xy$ at the point where $(x,y)=(1,2)$
\item the surface $yz=x^{2}$ at $(4,2,8)$
      \ {\itshape\small(Suggestion: first rewrite the surface as the graph of a function of
      two variables.)}
\end{parts}
\src{MIT 18.02, Problem Set 3; MA1521 L2, Q07.8, Q07.9}
\end{problem}
\work{6.8cm}

\begin{problem}
Consider the surface in $xyz$-space defined by
\[
\cos(\pi x)-x^{2}y+e^{xz}+yz=4 .
\]
Find equations for the tangent plane and the normal line to the surface at the point
$(0,1,2)$.
\src{Thomas 13.6, Exercise 5; MA1521 L2, Q10.1}
\end{problem}
\work{4.9cm}

\begin{problem}
\textbf{Tangent planes to a cone.}
\begin{parts}
\item Find an equation of the tangent plane to the cone $z=\sqrt{x^{2}+y^{2}}$ at a point
      $P_0:(x_0,y_0,z_0)$ on the cone, with $z_0>0$.
\item Write parametric equations for the ray from the origin passing through $P_0$, and use
      them to show that the ray lies on both the cone and the tangent plane at $P_0$.
\end{parts}
\src{MIT 18.02, Problem Set 3}
\end{problem}
\work{6.1cm}

\topichead{Differentials}

\begin{problem}
Find the differential ($dw$ or $dz$). Make the answer look as neat as possible.
\begin{parts}
\item $w=\ln(xyz)$
\item $w=xy^{2}z^{3}$
\item $\displaystyle z=\frac{x-y}{x+y}$
\end{parts}
\src{MIT 18.02, Problem Set 4; MA1521 L2, Q10.5}
\end{problem}
\work{4.4cm}

\begin{problem}
The following equations define $w$ implicitly as a function of the other variables. Find
$dw$ in terms of the other differentials by taking the differential of both sides and solving
algebraically for $dw$.
\begin{parts}
\item $z^{2}+xyw+yw^{3}=0$
\item $u^{2}+2v^{2}+3w^{2}=10$
\end{parts}
\src{MIT 18.02, Problem Set 4; MA1521 L2, Q10.6}
\end{problem}
\work{5.4cm}

\topichead{The Chain Rule}

\begin{problem}
In each of the following, find $dw/dt$ in two distinct ways: (i) use the Chain Rule, then
express the answer in terms of $t$; (ii) express the composite function in terms of $t$ and
differentiate.
\begin{parts}
\item $w=xyz$;\quad $x=t$, $y=t^{2}$, $z=t^{3}$
\item $w=x^{2}-y^{2}$;\quad $x=\cos t$, $y=\sin t$
\item $w=\ln\!\left(u^{2}+v^{2}\right)$;\quad $u=2\cos t$, $v=2\sin t$
\end{parts}
\src{MIT 18.02, Problem Set 4}
\end{problem}
\work{6.3cm}

\begin{problem}
Consider the function $w=2ye^{x}-\ln z$, where
$x=\ln\!\left(t^{2}+1\right)$, $y=\tan^{-1}t$, and $z=e^{t}$.
\begin{parts}
\item Find $dw/dt$ using the Chain Rule, expressing the answer in terms of $t$.
\item Find $dw/dt$ by first expressing $w$ in terms of $t$, and check that the two answers agree.
\end{parts}
\src{MA1521 L2, Q09.2}
\end{problem}
\work{5.8cm}

\begin{problem}
Suppose the temperature at any point in $xyz$-space is
$T(x,y,z)=e^{x}-yz+zx$, measured in ${}^{\circ}$C, where $x,y,z$ are measured in cm.
If the position of an insect at time $t$ seconds is
$x=t$, $y=2\left(1+t^{2}\right)$, $z=3+t$, find the rate of change of temperature
experienced by the insect at time $t=0$. Give the answer in ${}^{\circ}$C per second.
\src{MA1521 L2, Q09.3}
\end{problem}
\work{4.4cm}

\begin{problem}
Let $w=f(x,y)$, where all first- and second-order partial derivatives of $f$ are continuous,
and change to polar coordinates $x=r\cos\theta$, $y=r\sin\theta$.
\begin{parts}
\item Show that $\displaystyle (w_x)^{2}+(w_y)^{2}=(w_r)^{2}+\frac{1}{r^{2}}(w_\theta)^{2}$.
\item Show that
      $\displaystyle \pd{^{2}w}{r^{2}}=f_{xx}\cos^{2}\theta+2f_{xy}\cos\theta\sin\theta+f_{yy}\sin^{2}\theta$.
      {\itshape\small(First find $\pdi{w}{r}$ in terms of $f_x$, $f_y$, and $\theta$. Clairaut's
      Theorem may be used freely.)}
\end{parts}
\src{MIT 18.02, Problem Set 4; MA1521 L2, Q09.6}
\end{problem}
\work{7.1cm}

\begin{problem}
Let $w=f(x,y)$ and make the change of variables $x=u^{2}-v^{2}$, $y=2uv$. Show that
\[
(w_u)^{2}+(w_v)^{2}=4\left(u^{2}+v^{2}\right)\left[(w_x)^{2}+(w_y)^{2}\right].
\]
\src{MIT 18.02, Problem Set 4}
\end{problem}
\work{4.9cm}

\begin{problem}
\textbf{Partial differential equations from the Chain Rule.}
\begin{parts}
\item Let $w=f(y/x)$; that is, $w$ is the composite of $w=f(u)$ and $u=y/x$. Show that $w$
      satisfies the partial differential equation
      $\displaystyle x\,\pd{w}{x}+y\,\pd{w}{y}=0$.
\item Let $w=f\!\left(x^{2}-y^{2}\right)$; show that
      $\displaystyle y\,\pd{w}{x}+x\,\pd{w}{y}=0$.
\item Let $w=f(ax+by)$; show that
      $\displaystyle b\,\pd{w}{x}-a\,\pd{w}{y}=0$.
\end{parts}
\src{MIT 18.02, Problem Set 4}
\end{problem}
\work{6.3cm}

\topichead{Directional Derivatives and Gradient Vectors}

\begin{problem}
In each of the following, a function $f$, a point $P$, and a vector $\bA$ are given. Calculate
the gradient of $f$ at $P$, and the directional derivative $D_{\bu}f$ at $P$ in the direction
$\bu$ of the given vector $\bA$.
\begin{parts}
\item $f(x,y)=2xy-3y^{2}$;\quad $P(5,5)$,\quad $\bA=4\bi+3\bj$
\item $f(x,y,z)=xy+yz+zx$;\quad $P(1,-1,2)$,\quad $\bA=3\bi+6\bj-2\bk$
\item $f(x,y,z)=x^{2}+2y^{2}-3z^{2}$;\quad $P(1,1,1)$,\quad $\bA=\bi+\bj+\bk$
\item $f(u,v,w)=(u+2v+3w)^{2}$;\quad $P(1,-1,1)$,\quad $\bA=-2\bi+2\bj-\bk$
\end{parts}
\src{Thomas 13.5, Exercises 11, 15, 16; MIT 18.02, Problem Set 4}
\end{problem}
\work{6.8cm}

\begin{problem}
Consider the function $f(x,y)=2ye^{-x}+1$ and the point $P(0,-2)$.
\begin{parts}
\item Find the directional derivative of $f$ at $P$ in the direction of each of the vectors
      (i) $\tfrac{1}{\sqrt{2}}\bi+\tfrac{1}{\sqrt{2}}\bj$; \quad (ii) $12\bi+5\bj$.
\item Find the direction that gives the largest possible directional derivative of $f$ at $P$,
      and the value of that derivative.
\item Find the direction along which $f$ decreases most rapidly at $P$.
\item Find the directions along which $f$ exhibits no change at $P$.
\end{parts}
\src{MA1521 L2, Q09.10}
\end{problem}
\work{6.8cm}

\begin{problem}
Let $a$ be a positive number and consider
$f(x,y)=\left(x^{2}+a\right)e^{2x+3y^{2}}$.
It is given that the directional derivative of $f$ at the point $(-6,2)$ in the direction of
the vector $\bi+2\bj$ is $216\sqrt{5}$. Find the value of $a$.
\src{MA1521 L2, Q09.9}
\end{problem}
\work{5.3cm}

\begin{problem}
Let $C$ be the curve of intersection of the two surfaces
$x^{3}+2xy+yz=7$ and $3x^{2}-yz=1$.
Find parametric equations of the tangent line to $C$ at $(1,2,1)$.
\par\vspace{3pt}
{\itshape\small Suggestion. For each of the two surfaces, first find a normal vector at the
point $(1,2,1)$.}
\src{Thomas 13.6, Exercises 15--20; MA1521 L2, Q10.2}
\end{problem}
\work{5.4cm}

\topichead{Linearization and Quadratic Approximation}

\begin{problem}
\textbf{Linearization and linear approximation.}
\begin{parts}
\item Find the linearization of $f(x,y)=e^{2x+y}$ at $(0,0)$ and at $(1,2)$. Give each answer
      in the form $ax+by+c$.
\item Consider $f(x,y,z)=xz+\sin(xyz)$. Using linear approximation, estimate the change in the
      value of $f$ when $(x,y,z)$ changes from
      $\left(\tfrac{1}{2},\pi,\tfrac{1}{3}\right)$ to
      $\left(\tfrac{1}{2}+0.03,\ \pi+0.04,\ \tfrac{1}{3}-0.02\right)$.
\end{parts}
\src{MA1521 L2, Q10.4, Q10.7}
\end{problem}
\work{6.8cm}

\begin{problem}
Consider the function $f(x,y)=e^{x+4y}\cos x$.
\begin{parts}
\item Find the second-degree Taylor polynomial of $f$ at $(0,0)$.
\item Using part (a), find an approximate value of $f(0.2,\,0.1)$.
\end{parts}
\src{Thomas 13.9; MA1521 L2, Q10.9}
\end{problem}
\work{6.5cm}


%%%%%%%%%%%%%%%%%%%%%%%%%%%%%%%%%%%%%%%%%%%%%%%%%%%%%%%%%%%%%%%%%%% PART D %%%
\parthead{D}{Extreme Values, Saddle Points, and Lagrange Multipliers}%
{Thomas 13.7--13.9 \ \textbullet\ \ MIT 18.02 Lectures 9, 10, 13}%
{14 problems \ \textbar\ \ 60 points}

\topichead{Finding Local Extrema}

\begin{problem}
Find all the local maxima, local minima, and saddle points of the following functions, and
give the corresponding values of the function.
\begin{parts}
\item $f(x,y)=x^{2}+xy+y^{2}+3x-3y+4$
\item $f(x,y)=x^{3}-y^{3}-2xy+6$
\item $f(x,y)=4xy-x^{4}-y^{4}$
\item $f(x,y)=2\ln x+\ln y-4x-y$ \quad ($x>0$, $y>0$)
\end{parts}
\src{Thomas 13.7, Exercises 1, 13, 19, 29}
\end{problem}
\work{7.8cm}

\begin{problem}
Determine all the critical points of
\[
f(x,y)=2xy-x^{2}y-xy^{2},
\]
and classify each as a local maximum point, a local minimum point, or a saddle point of $f$.
Evaluate also the local maximum and local minimum values of $f$, if any.
\src{MA1521 L2, Q08.2}
\end{problem}
\work{7.1cm}

\begin{problem}
Determine all the critical points of
\[
f(x,y)=\ln\!\left(xy^{2}\right)-xy-2y
\]
over the region $x>0$, $y>0$ in the $xy$-plane, and classify each one. Evaluate also the local
maximum and local minimum values of $f$, if any.
\src{MA1521 L2, Q08.3}
\end{problem}
\work{6.8cm}

\topichead{Finding Absolute Extrema}

\begin{problem}
Find the maximum and minimum points of
\[
f(x,y)=x^{2}-2xy+2y^{2}-2y
\]
over the rectangular region $R:\ -1\le x\le 2$, $0\le y\le 2$ in the $xy$-plane. Hence find the
absolute maximum and absolute minimum values of $f$ over $R$.
\src{MIT 18.02, Problem Set 3; MA1521 L2, Q08.6}
\end{problem}
\work{8.2cm}

\begin{problem}
\textbf{Constrained minimum distance.}
\begin{parts}
\item Consider the surface $y^{2}-xz=4$ in $xyz$-space. Find the point or points on the
      surface closest to the origin.
      {\itshape\small(Suggestion: first formulate the question as a problem of minimizing a
      function of the variables $x$ and $z$.)}
\item Find the minimum distance from the point $(2,-1,1)$ to the plane $x+y-z=2$.
\end{parts}
\src{Thomas 13.7, Exercise 52; MA1521 L2, Q08.4}
\end{problem}
\work{6.5cm}

\begin{problem}
\textbf{The method of least squares.}
\begin{parts}
\item Using the method of least squares, find the line $y=ax+b$ that best fits the three data
      points $(-1,1)$, $(1,2)$, $(3,2)$.
\item Using the method of least squares, find the quadratic polynomial of the form
      $y=ax^{2}+bx$ that best fits the three data points $(-1,1)$, $(1,1)$, $(2,0)$.
\end{parts}
\src{MIT 18.02, Problem Set 3; MA1521 L2, Q08.7, Q08.8}
\end{problem}
\work{7.1cm}

\begin{problem}
A rectangular box with its top left open has its five faces made from different materials.
The front face and the back face are made from material costing $10$ dollars per square metre;
the two lateral faces from material costing $5$ dollars per square metre; and the bottom face
from material costing $30$ dollars per square metre. Find the minimum material cost of the box
if its volume is $6$ cubic metres, and give the dimensions (length, width, and height) of such
a box.
\par\vspace{3pt}
{\itshape\small Do this problem twice: first by eliminating a variable and using the Second
Derivative Test, and then again by the method of Lagrange multipliers. Compare the work
involved.}
\src{MIT 18.02, Problem Set 3; MA1521 L2, Q08.5 and Q11.4}
\end{problem}
\work{8.8cm}

\topichead{The Method of Lagrange Multipliers}

\begin{problem}
\textbf{Extrema on a circle.}
\begin{parts}
\item Use the method of Lagrange multipliers to find the maximum value of $f(x,y)=xy$
      subject to the constraint $(x-2)^{2}+y^{2}=4$.
\item Use the method of Lagrange multipliers to find the maximum value of $f(x,y)=xy$
      subject to the constraints $(x-2)^{2}+y^{2}=4$ \emph{and} $x\le 2$.
      {\itshape\small(This is part (a) with one extra constraint, so the work there is useful.)}
\end{parts}
\src{MA1521 L2, Q11.1, Q11.2}
\end{problem}
\work{7.8cm}

\begin{problem}
\textbf{Nearest and farthest points.}
\begin{parts}
\item Find the point or points on the sphere $x^{2}+y^{2}+z^{2}=1$ furthest from the point
      $(2,1,2)$.
\item Find the point on the plane $x+2y+3z=13$ closest to the point $(1,1,1)$.
\end{parts}
\src{Thomas 13.8, Exercise 17; MA1521 L2, Q11.3}
\end{problem}
\work{6.8cm}

\begin{problem}
A rectangular box is placed in the first octant so that one corner $Q$ is at the origin and the
three sides adjacent to $Q$ lie in the coordinate planes. The corner $P$ diagonally opposite
to $Q$ lies on the ellipsoid
\[
\frac{x^{2}}{3}+\frac{y^{2}}{9}+\frac{z^{2}}{4}=1 .
\]
Use the method of Lagrange multipliers to determine the largest possible volume of the box, and
say how you know that your answer gives a maximum.
\src{MIT 18.02, Problem Set 5; MA1521 L2, Q12.1}
\end{problem}
\work{7.8cm}

\begin{problem}
Let $f(x,y)=x^{2}-y^{2}$ be defined on the disk $D=\left\{(x,y):x^{2}+y^{2}\le1\right\}$.
\begin{parts}
\item Find the critical points of $f$ and classify them.
\item Find the maximum and minimum values of $f$ on $D$.
      {\itshape\small(Lagrange multipliers are used in part (b).)}
\end{parts}
\src{MA1521 L2, Q12.2}
\end{problem}
\work{7.1cm}

\begin{problem}
\textbf{Extrema on a curve of intersection.}
\begin{parts}
\item Find the extreme values of $f(x,y,z)=xy+z^{2}$ on the circle in which the plane
      $y-x=0$ intersects the sphere $x^{2}+y^{2}+z^{2}=4$.
\item Find the maximum value of $w=xyz$ on the line of intersection of the two planes
      $x+y+z=40$ and $x+y-z=0$. Give a geometric argument supporting the claim that you have
      found a maximum and not a minimum.
\end{parts}
\src{Thomas 13.8, Exercises 42, 43}
\end{problem}
\work{7.8cm}

\begin{problem}
\textbf{Three applications.}
\begin{parts}
\item The temperature at a point $(x,y)$ on a metal plate is $T(x,y)=4x^{2}-4xy+y^{2}$. An ant
      on the plate walks around the circle of radius $5$ centred at the origin. What are the
      highest and lowest temperatures the ant encounters?
\item Find the dimensions of the closed right circular cylindrical can of smallest surface area
      whose volume is $16\pi\ \text{cm}^{3}$.
\item Suppose the utility of amounts $x$ and $y$ of two capital goods is $U(x,y)=xy+2x$, and
      that the budget constraint simplifies to $2x+y=30$. Find the maximum value of $U$ and the
      corresponding values of $x$ and $y$.
\end{parts}
\src{Thomas 13.8, Exercises 9, 15, 33}
\end{problem}
\work{8.2cm}

\begin{problem}
\textbf{The condition $\grad f=\lambda\grad g$ is not sufficient.}\quad
Even though $\grad f=\lambda\grad g$ is a necessary condition for an extreme value of
$f(x,y)$ subject to $g(x,y)=0$ and $\grad g\neq\mathbf{0}$, it does not by itself guarantee
that an extreme value exists. Try using Lagrange multipliers to find a maximum value of
$f(x,y)=x+y$ subject to $xy=16$. The method identifies $(4,4)$ and $(-4,-4)$ as candidates.
Explain carefully why $x+y$ nevertheless has \emph{no} maximum value on the hyperbola
$xy=16$, and say which of the two candidates is a genuine local extremum and of what kind.
\src{Thomas 13.8, Exercise 45}
\end{problem}
\work{6.8cm}


%%%%%%%%%%%%%%%%%%%%%%%%%%%%%%%%%%%%%%%%%%%%%%%%%%%%%%%%%%%%%%%%%%% PART E %%%
\parthead{E}{Midterm Simulation}%
{Thomas 11--13 \ \textbullet\ \ 90 minutes, no calculator, closed book}%
{6 problems \ \textbar\ \ 100 points}

\panel{INSTRUCTIONS}{\fontsize{10}{12.5}\selectfont
Work the six problems below in one uninterrupted sitting of \textbf{90 minutes}. Show all
reasoning; an unsupported answer earns no credit. Leave answers in exact form ($\sqrt{2}$,
$\ln 3$, $\pi/4$, and so on). The point values sum to $100$.}

\begin{problem}
\hfill\pts{15 points}\par\vspace{2pt}
Let $A=(1,0,2)$, $B=(2,1,0)$, and $C=(0,3,1)$.
\begin{parts}
\item Find the area of the triangle $ABC$.
\item Find an equation of the plane containing $A$, $B$, and $C$.
\item Find the distance from the origin to that plane.
\end{parts}
\end{problem}
\work{7.6cm}

\begin{problem}
\hfill\pts{13 points}\par\vspace{2pt}
Consider the plane curve
\[
\br(t)=(\cos t+t\sin t)\,\bi+(\sin t-t\cos t)\,\bj ,\qquad t\ge0 .
\]
\begin{parts}
\item Find $\bv(t)$ and the speed $|\bv(t)|$, simplifying as far as possible.
\item Find the arc length parameter $s(t)$ measured from the point where $t=0$, and hence the
      length of the portion of the curve with $0\le t\le\pi/2$.
\item Find the unit tangent vector $\bT(t)$ for $t>0$.
\end{parts}
\end{problem}
\work{7.6cm}

\begin{problem}
\hfill\pts{18 points}\par\vspace{2pt}
Let $f(x,y,z)=x^{2}y+yz^{2}$ and let $P=(1,2,-1)$.
\begin{parts}
\item Find $\grad f$ at $P$.
\item Find the directional derivative of $f$ at $P$ in the direction of $2\bi-\bj+2\bk$.
\item In what direction does $f$ increase most rapidly at $P$, and what is that rate of increase?
\item A particle moves so that its position at time $t$ is
      $\br(t)=t\,\bi+(t+1)\,\bj+\left(t^{2}-2\right)\bk$. Verify that the particle is at $P$ when
      $t=1$, and find the rate of change of $f$ experienced by the particle at that instant.
\end{parts}
\end{problem}
\work{8.5cm}

\begin{problem}
\hfill\pts{16 points}\par\vspace{2pt}
Consider the ellipsoid $x^{2}+2y^{2}+3z^{2}=21$ and the point $P_0(1,2,2)$ on it.
\begin{parts}
\item Find an equation of the tangent plane and parametric equations of the normal line to the
      ellipsoid at $P_0$.
\item Regarding $z$ as a function of $x$ and $y$ near $P_0$, use implicit differentiation to find
      $\pdi{z}{x}$ and $\pdi{z}{y}$ at $P_0$, and check that your answers agree with part (a).
\item Find all points on the ellipsoid at which the tangent plane is parallel to the plane
      $x+4y+6z=0$.
\end{parts}
\end{problem}
\work{8.5cm}

\begin{problem}
\hfill\pts{18 points}\par\vspace{2pt}
Find and classify all the critical points of
\[
f(x,y)=x^{3}+3xy^{2}-15x-12y ,
\]
using the Second Derivative Test. State the local maximum and local minimum values of $f$.
\end{problem}
\work{9.0cm}

\begin{problem}
\hfill\pts{20 points}\par\vspace{2pt}
\begin{parts}
\item Use the method of Lagrange multipliers to find the maximum and minimum values of
      $f(x,y,z)=x-2y+5z$ on the sphere $x^{2}+y^{2}+z^{2}=30$.
\item Find the dimensions of the closed rectangular box of maximum volume that can be inscribed
      in the sphere $x^{2}+y^{2}+z^{2}=4$, and state that maximum volume.
\end{parts}
\end{problem}
\work{9.0cm}


%%%%%%%%%%%%%%%%%%%%%%%%%%%%%%%%%%%%%%%%%%%%%%%%%%%%%%%%%%%%%%%%%% ANSWERS %%%
\clearpage
\markboth{Answers}{Answers}
\noindent{\color{ThomasBlue}\rule{\linewidth}{2.6pt}}\par\vspace{2pt}
\noindent{\color{RuleGray}\rule{\linewidth}{0.5pt}}\par\vspace{7pt}
\noindent{\sffamily\bfseries\fontsize{21}{24}\selectfont\color{ThomasBlue}Answers}%
\hfill{\sffamily\fontsize{9}{11}\selectfont\color{SoftGray}Exercise Set M}\par\vspace{4pt}
\noindent{\itshape\fontsize{10}{12}\selectfont\color{SoftGray}
Consult these only after a complete attempt. ``Show that'' problems are answered with the
key step rather than the full argument.}\par\vspace{5pt}
\noindent{\color{ThomasBlue}\rule{\linewidth}{1.1pt}}\par\vspace{10pt}

{\fontsize{8.8}{11}\selectfont
\setlength{\columnsep}{22pt}
\setlength{\parskip}{3.5pt}
\begin{multicols}{2}
\raggedright

\textbf{Part A}\par\vspace{2pt}

\textbf{1.} (a) $c=\pm\sqrt{6}/4$.\quad (b) $(0,4,-3)$.

\textbf{2.} (b) $\overrightarrow{OX}=s\,\overrightarrow{OP}+r\,\overrightarrow{OQ}$.

\textbf{3.} (a) $\pi/4$ and $\pi/3$.\quad (b) The dot product is $c-4$; (i) $c=4$, (ii) $c>4$.

\textbf{4.} $\bC=-4\bi+2\bj+2\bk$, $\bD=4\bi+4\bj+4\bk$.

\textbf{5.} (a) $-34$ by both expansions.\quad
(b) (i) $3\bi+3\bj+3\bk$; (ii) $3\bi-\bj+2\bk$.

\textbf{6.} (a) $\overrightarrow{AB}\times\overrightarrow{AC}=26\bi+18\bj+6\bk$; parallelogram
$2\sqrt{259}\approx32.19$, triangle $\sqrt{259}$.\quad (b) $11$.\quad
(c) Yes --- a zero determinant means the three edge vectors are linearly dependent, so all four
points lie in a common plane.

\textbf{7.} (a) $x+2y-2z=4$.\quad (b) $x-y-z=0$.\quad
(c) $\dfrac{x}{a}+\dfrac{y}{b}+\dfrac{z}{c}=1$.\quad (d) $x+y=1$.

\textbf{8.} (a) $x=1+2t$, $y=-t$, $z=-1+3t$.\quad
(b) $x=2+t$, $y=-1-t$, $z=-1+2t$.\quad
(c) $x=1+at$, $y=1+bt$, $z=1+(a+2b)t$ for any $(a,b)\neq(0,0)$.\quad
(d) Direction $\langle3,-4,4\rangle\times\langle3,-1,5\rangle=\langle-16,-3,9\rangle$:
$x=4-16t$, $y=2-3t$, $z=1+9t$.

\textbf{9.} (a) $(4,-1,4)$.\quad (b) $(0,-1,0)$.\quad (c) $x-2y+4z+1=0$.

\vspace{4pt}
\textbf{Part B}\par\vspace{2pt}

\textbf{10.} (a) $\bv=e^{t}\bi-e^{-t}\bj$, speed $\sqrt{e^{2t}+e^{-2t}}$,
$\bT=\bv/|\bv|$, $\ba=e^{t}\bi+e^{-t}\bj$.\quad
(b) $\bv=2t\bi+3t^{2}\bj$, speed $|t|\sqrt{4+9t^{2}}$,
$\bT=(2\bi+3t\bj)/\sqrt{4+9t^{2}}$ for $t>0$, $\ba=2\bi+6t\bj$.\quad
(c) $\bv=t(-4\bi+2\bj+4\bk)$, speed $6|t|$, $\bT=\tfrac{1}{3}(-2\bi+\bj+2\bk)$ for $t>0$,
$\ba=-4\bi+2\bj+4\bk$ (constant).

\textbf{11.} (a) $\bv'=(6t+2)\bi+\pi^{2}\cos(\pi t)\bj+e^{t}\bk$;\ $\bv'(2)=14\bi+\pi^{2}\bj+e^{2}\bk$.\quad
(b) $\left(t^{3}+t^{2}\right)\bi-\cos(\pi t)\bj+e^{t}\bk+\bC$.\quad
(c) $2\bi+2\bj+(e-1)\bk$.

\textbf{12.} $\br(1)=\mathbf{0}$ and $\br'(1)=\bi+\tfrac14\bj+\bk$, so
$x=t$, $y=t/4$, $z=t$.

\textbf{13.} $5\bi+5\bj+\dfrac{5\pi}{4}\bk$.

\textbf{14.} (a) No: $t=1+2t$ forces $t=-1$, and then $t^{2}=1\neq-5$.\quad
(b) Yes. Solving $s=1+2t$, $s^{2}=1+6t$, $s^{3}=1+14t$ gives $s=1$ ($t=0$) and $s=2$
($t=\tfrac12$); the paths meet at $(1,1,1)$ and $(2,4,8)$.

\textbf{15.} (a) $|\bv|=3$, $\bT=\tfrac{1}{3}\left(-2\sin t\,\bi+2\cos t\,\bj+\sqrt{5}\,\bk\right)$,
length $3\pi$.\quad
(b) $|\bv|=4\left|\sin(t/2)\right|$, length $32$.\quad
(c) $\sqrt{3}+\ln\!\left(2+\sqrt{3}\right)\approx3.049$.

\vspace{4pt}
\textbf{Part C}\par\vspace{2pt}

\textbf{16.} (a) $4/5$ and $3$.\quad
(b) $\left\{(x,y):0\le x^{2}+y^{2}<9\ \text{and}\ x^{2}+y^{2}\neq8\right\}$: the open disk of
radius $3$ with the circle of radius $2\sqrt{2}$ removed.\quad
(c) For $0\le c<4$, the sphere $x^{2}+y^{2}+z^{2}=16-c^{2}$ centred at the origin; for $c=4$,
the origin alone.

\textbf{17.} (a) (i) parallel lines $x+y=1-c$, graph a plane; (ii) circles $x^{2}+y^{2}=c^{2}$,
$c\ge0$, graph a cone; (iii) circles of radius $\sqrt{c}$, graph a paraboloid opening upward;
(iv) circles $x^{2}+y^{2}=1-c$, graph a paraboloid opening downward.\quad
(b) Lines through the origin (origin and the line $y=x$ deleted):
$y=\dfrac{k-1}{k+1}\,x$, i.e. $y=3x$, $y=-x$, $y=0$, $y=\tfrac13x$ for $k=-2,0,1,2$.

\textbf{18.} (a) $1$.\quad (b) $\dfrac14$ (multiply by
$\sqrt{2x-y}+2$).\quad (c) Does not exist: the limit is $0$ along $y=0$ but $1$ along
$y=x^{2}$.\quad (d) The limit is $0\neq f(0,0)=1$, so $f$ is not continuous at the origin.

\textbf{19.} (a) $12x^{2}-6xy^{2}+2$;\ $-6x^{2}y$.\quad
(b) $3\cos(3x+2y)$;\ $2\cos(3x+2y)$.\quad
(c) $\dfrac{4y}{4x+3y}$;\ $\ln(4x+3y)+\dfrac{3y}{4x+3y}$; at $(2,1)$: $\tfrac{4}{11}$ and
$\ln11+\tfrac{3}{11}$.\quad
(d) $2xz$;\ $-2z^{3}$;\ $x^{2}-6yz^{2}$.\quad
(e) $3x^{2}yz^{3}$;\ $2yz+x^{3}z^{3}$;\ $y^{2}+3x^{3}yz^{2}$; at $(1,2,4)$: $384$, $80$, $100$.

\textbf{20.} (b) $f_{xy}=ax+6y$ and $f_{yx}=2x+6y$, so $a=2$ and $f=x^{2}y+3xy^{2}+C$.\quad
(c) $f_{xy}=Ax$ and $f_{yx}=6x$, so $A=6$ and $f=x^{3}+3x^{2}y+2y^{3}+C$.

\textbf{21.} (b) $\pd{z}{x}=\dfrac{3x^{2}+6z^{3}}{59-18xz^{2}}$,\quad
$\pd{z}{y}=\dfrac{6y^{2}}{59-18xz^{2}}$.

\textbf{22.} (a) $z=x+2y-2$.\quad (b) $w=-4x+4y$.\quad (c) $z=8x+3y-8$.\quad
(d) $4x-4y-z=0$.

\textbf{23.} $\grad F=\langle2,2,1\rangle$ at $(0,1,2)$; tangent plane $2x+2y+z=4$, normal line
$x=2t$, $y=1+2t$, $z=2+t$.

\textbf{24.} (a) $z_{0}z=x_{0}x+y_{0}y$.\quad
(b) The ray is $(tx_{0},ty_{0},tz_{0})$, $t\ge0$; substituting gives $t z_{0}=t z_{0}$ on the
cone and $t z_{0}^{2}=t\!\left(x_{0}^{2}+y_{0}^{2}\right)=tz_{0}^{2}$ on the plane. (A cone is
tangent to each of its tangent planes along a whole ruling.)

\textbf{25.} (a) $\dfrac{dx}{x}+\dfrac{dy}{y}+\dfrac{dz}{z}$.\quad
(b) $xy^{2}z^{3}\!\left(\dfrac{dx}{x}+\dfrac{2\,dy}{y}+\dfrac{3\,dz}{z}\right)$.\quad
(c) $\dfrac{2y\,dx-2x\,dy}{(x+y)^{2}}$.

\textbf{26.} (a) $dw=-\dfrac{yw\,dx+\left(xw+w^{3}\right)dy+2z\,dz}{xy+3yw^{2}}$.\quad
(b) $dw=-\dfrac{u\,du+2v\,dv}{3w}$.

\textbf{27.} (a) $6t^{5}$.\quad (b) $-2\sin2t$.\quad (c) $0$ (here $u^{2}+v^{2}\equiv4$).

\textbf{28.} $\dd{w}{t}=4t\tan^{-1}t+1$ by either method.

\textbf{29.} $\grad T=\langle4,-3,-2\rangle$ at $(0,2,3)$ and
$\br'(0)=\langle1,0,1\rangle$, so $dT/dt=2\ {}^{\circ}\text{C/s}$.

\textbf{30.} (a) $w_r=f_x\cos\theta+f_y\sin\theta$ and
$w_\theta=r\left(-f_x\sin\theta+f_y\cos\theta\right)$; squaring and adding gives the identity.\quad
(b) Differentiate $w_r$ again with respect to $r$, using the Chain Rule on $f_x$ and $f_y$ and
$f_{xy}=f_{yx}$.

\textbf{31.} $w_u=2u\,w_x+2v\,w_y$ and $w_v=-2v\,w_x+2u\,w_y$; the cross terms cancel when
the squares are added.

\textbf{32.} (a) $w_x=-\dfrac{y}{x^{2}}f'$, $w_y=\dfrac1x f'$.\quad
(b) $w_x=2xf'$, $w_y=-2yf'$.\quad (c) $w_x=af'$, $w_y=bf'$.

\textbf{33.} (a) $\grad f=\langle10,-20\rangle$, $D_{\bu}f=-4$.\quad
(b) $\grad f=\langle1,3,0\rangle$, $D_{\bu}f=3$.\quad
(c) $\grad f=\langle2,4,-6\rangle$, $D_{\bu}f=0$.\quad
(d) $\grad f=\langle4,8,12\rangle$, $D_{\bu}f=-\tfrac43$.

\textbf{34.} $\grad f(P)=\langle4,2\rangle$.
(a) (i) $3\sqrt{2}$; (ii) $\tfrac{58}{13}$.\quad
(b) $\tfrac{1}{\sqrt5}\langle2,1\rangle$, with value $|\grad f|=2\sqrt5$.\quad
(c) $-\tfrac{1}{\sqrt5}\langle2,1\rangle$, with value $-2\sqrt5$.\quad
(d) $\pm\tfrac{1}{\sqrt5}\langle1,-2\rangle$.

\textbf{35.} $\grad f(-6,2)=\langle60+2a,\ 432+12a\rangle$ and
$D_{\bu}f=(924+26a)/\sqrt5=216\sqrt5$, so $a=6$.

\textbf{36.} $\grad F=\langle7,3,2\rangle$, $\grad G=\langle6,-1,-2\rangle$, and
$\grad F\times\grad G=\langle-4,26,-25\rangle$:
$x=1-4t$, $y=2+26t$, $z=1-25t$.

\textbf{37.} (a) $L(x,y)=1+2x+y$ at $(0,0)$; \
$L(x,y)=e^{4}(2x+y-3)$ at $(1,2)$.\quad
(b) $df=\left(\tfrac13+\tfrac{\sqrt3\pi}{6}\right)(0.03)+\tfrac{\sqrt3}{12}(0.04)
-\left(\tfrac12+\tfrac{\sqrt3\pi}{4}\right)(0.02)=\dfrac{\sqrt3}{300}\approx0.0058$.

\textbf{38.} (a) $1+x+4y+4xy+8y^{2}$.\quad
(b) $1.76$ (the true value is $1.7858\ldots$).

\vspace{4pt}
\textbf{Part D}\par\vspace{2pt}

\textbf{39.} (a) Local minimum $-5$ at $(-3,3)$ ($D=3$, $f_{xx}=2$).\quad
(b) Saddle at $(0,0)$ ($D=-4$); local maximum $\tfrac{170}{27}$ at
$\left(-\tfrac23,\tfrac23\right)$ ($D=12$, $f_{xx}=-4$).\quad
(c) Saddle at $(0,0)$; local maxima of value $2$ at $(1,1)$ and $(-1,-1)$.\quad
(d) Local maximum $-3-2\ln2$ at $\left(\tfrac12,1\right)$ ($D=8$, $f_{xx}=-8$).

\textbf{40.} Critical points $(0,0)$, $(2,0)$, $(0,2)$, $\left(\tfrac23,\tfrac23\right)$;
$D=4xy-(2-2x-2y)^{2}$. The first three are saddle points ($D=-4$); the last is a local maximum
with value $\tfrac{8}{27}$ ($D=\tfrac43$, $f_{xx}=-\tfrac43$). There is no local minimum.

\textbf{41.} The only critical point is $\left(2,\tfrac12\right)$, with $D=1>0$ and
$f_{xx}=-\tfrac14<0$: a local maximum of value $-2-\ln2$. There is no local minimum and no
saddle point.

\textbf{42.} Interior critical point $(1,1)$ with $f=-1$. On the boundary the candidates are
$f(-1,0)=1$, $f(-1,2)=9$, $f(2,0)=4$, $f(2,2)=0$, $f\!\left(2,\tfrac32\right)=-\tfrac12$,
$f(0,0)=0$. Absolute maximum $9$ at $(-1,2)$; absolute minimum $-1$ at $(1,1)$.

\textbf{43.} (a) Minimise $x^{2}+y^{2}+z^{2}=x^{2}+xz+z^{2}+4$; the only critical point is
$x=z=0$, giving $(0,\pm2,0)$ at distance $2$.\quad
(b) $\dfrac{|2-1-1-2|}{\sqrt3}=\dfrac{2\sqrt3}{3}$.

\textbf{44.} (a) $y=\tfrac14x+\tfrac{17}{12}$.\quad
(b) $y=\tfrac{3}{11}x^{2}-\tfrac{4}{11}x$.

\textbf{45.} With width $x$, depth $y$, height $z$: $C=20xz+10yz+30xy$ subject to $xyz=6$.
Eliminating $z$ gives $C=\dfrac{60}{x}+\dfrac{120}{y}+30xy$, whose only critical point is
$x=1$, $y=2$, so $z=3$. Minimum cost $\$180$, with dimensions
$1\ \text{m}\times2\ \text{m}\times3\ \text{m}$ (each of the three cost terms equals $\$60$).

\textbf{46.} (a) The equations give $y^{2}=x^{2}-2x$, hence $x=0$ or $x=3$; the maximum is
$3\sqrt3$ at $\left(3,\sqrt3\right)$ (and the minimum is $-3\sqrt3$ at
$\left(3,-\sqrt3\right)$).\quad
(b) The candidate $x=3$ is excluded, so the maximum occurs on the new boundary $x=2$: the
maximum is $4$, at $(2,2)$.

\textbf{47.} (a) $\left(-\tfrac23,-\tfrac13,-\tfrac23\right)$, at distance $4$.\quad
(b) $\left(\tfrac32,2,\tfrac52\right)$.

\textbf{48.} At the maximum $\dfrac{x^{2}}{3}=\dfrac{y^{2}}{9}=\dfrac{z^{2}}{4}=\dfrac13$, so
$x=1$, $y=\sqrt3$, $z=\tfrac{2}{\sqrt3}$ and $V_{\max}=2$. The constraint surface is closed and
bounded and $V$ is continuous, so the extreme values are attained; $V=0$ on the boundary of the
first octant, so the interior critical point is the maximum.

\textbf{49.} (a) $(0,0)$ is the only critical point, and $D=-4<0$: a saddle point, $f(0,0)=0$.\quad
(b) On $x^{2}+y^{2}=1$: maximum $1$ at $(\pm1,0)$, minimum $-1$ at $(0,\pm1)$.

\textbf{50.} (a) On $y=x$ the constraint becomes $2x^{2}+z^{2}=4$ and $f=4-x^{2}$: maximum $4$
at $(0,0,\pm2)$, minimum $2$ at $\pm\left(\sqrt2,\sqrt2,0\right)$.\quad
(b) The planes meet where $x+y=20$, $z=20$, so $w=20xy$ with $x+y=20$; the maximum is
$w=2000$ at $(10,10,20)$. Along the line $w=20x(20-x)\to-\infty$ in both directions, so the
critical point must be a maximum.

\textbf{51.} (a) $T=(2x-y)^{2}$; on $x=5\cos t$, $y=5\sin t$ this is
$25(2\cos t-\sin t)^{2}$, so the hottest temperature is $125^{\circ}$ and the coldest is
$0^{\circ}$.\quad
(b) $r=2$ cm, $h=4$ cm.\quad
(c) $U_{\max}=128$ at $x=8$, $y=14$.

\textbf{52.} The method gives $x=y$, so $(4,4)$ and $(-4,-4)$. On the branch in the first
quadrant $f=x+16/x\to+\infty$ as $x\to0^{+}$ and as $x\to+\infty$, so $(4,4)$ is a
\emph{minimum} there ($f=8$) and $f$ has no maximum on $xy=16$. On the third-quadrant branch
$(-4,-4)$ gives a maximum, $f=-8$. The constraint set is unbounded, which is exactly what
$\grad f=\lambda\grad g$ cannot detect.

\vspace{4pt}
\textbf{Part E}\par\vspace{2pt}

\textbf{53.} $\overrightarrow{AB}\times\overrightarrow{AC}=\langle5,3,4\rangle$.
(a) $\tfrac{5\sqrt2}{2}$.\quad (b) $5x+3y+4z=13$.\quad
(c) $\dfrac{13}{5\sqrt2}=\dfrac{13\sqrt2}{10}$.

\textbf{54.} (a) $\bv=t\cos t\,\bi+t\sin t\,\bj$, $|\bv|=t$ for $t\ge0$.\quad
(b) $s(t)=t^{2}/2$; the length is $\pi^{2}/8$.\quad
(c) $\bT=\cos t\,\bi+\sin t\,\bj$. (This curve is the involute of the unit circle.)

\textbf{55.} (a) $\grad f(P)=\langle4,2,-4\rangle$.\quad
(b) $-\tfrac23$.\quad
(c) In the direction $\tfrac13\langle2,1,-2\rangle$, at the rate $|\grad f|=6$.\quad
(d) $\br(1)=\langle1,2,-1\rangle=P$ and $\br'(1)=\langle1,1,2\rangle$, so
$df/dt=4+2-8=-2$.

\textbf{56.} $\grad F(P_0)=\langle2,8,12\rangle$.
(a) Tangent plane $x+4y+6z=21$; normal line $x=1+t$, $y=2+4t$, $z=2+6t$.\quad
(b) $\pdi{z}{x}=-x/(3z)=-\tfrac16$ and $\pdi{z}{y}=-2y/(3z)=-\tfrac23$, matching
$z=\tfrac16(21-x-4y)$.\quad
(c) $\grad F\parallel\langle1,4,6\rangle$ forces $(x,y,z)=\pm(1,2,2)$, so the points are
$(1,2,2)$ and $(-1,-2,-2)$.

\textbf{57.} $f_x=0$ and $f_y=0$ give $x^{2}+y^{2}=5$ and $xy=2$, so the critical points are
$(2,1)$, $(1,2)$, $(-1,-2)$, $(-2,-1)$. Here $D=36\!\left(x^{2}-y^{2}\right)$.
Local minimum $-28$ at $(2,1)$; local maximum $28$ at $(-2,-1)$; saddle points at $(1,2)$ and
$(-1,-2)$.

\textbf{58.} (a) $\langle1,-2,5\rangle=2\lambda\langle x,y,z\rangle$ gives
$(x,y,z)=\pm(1,-2,5)$: maximum $30$ at $(1,-2,5)$, minimum $-30$ at $(-1,2,-5)$.\quad
(b) A cube of edge $\dfrac{4}{\sqrt3}$, of volume $\dfrac{64}{3\sqrt3}=\dfrac{64\sqrt3}{9}$.

\end{multicols}}

\vspace{10pt}
\noindent{\color{RuleGray}\rule{\linewidth}{0.5pt}}\par\vspace{5pt}
\noindent{\sffamily\fontsize{8}{10}\selectfont\color{SoftGray}
Problems adapted from \emph{Thomas' Calculus: Early Transcendentals}, 15th edition
(Hass, Heil, Weir), \S\S11.1--11.5, 12.1--12.3, 13.1--13.9; from the MIT 18.02 problem sets
distributed with the Lecture Group 1 lesson notes; and from the MA1521 Lecture Group 2 tutorial
questions, Lessons 02--12.}

\end{document}
